% !TEX root =./ECE444_Practice_main.tex
%
% L23 practice set (Array Factor and Pattern Lab), in two parts that follow the
% lab's two parts. Part A (objective 3.2; the former L21 Array Factor Lab set,
% merged 2026-10-11): reading HPBW, FNBW and sidelobe level off a PHASER beam
% sweep, predicting a reduced aperture before measuring it, and reconciling a
% measured sidelobe against the array-factor value. Part B (objective 3.6; the
% former L23 Antenna Pattern Lab set): reading a hand-rotated pattern trace:
% dBc extraction, why a time axis cannot be calibrated to angle, scan loss and
% sidelobe asymmetry, far-field radius, and what a chamber removes.

\fullwidth{\textbf{Part A -- Learning Objective 3.2: } I can derive the array factor for an arbitrary linear
array and apply pattern multiplication.
\\[4pt]\normalfont\small Show your work. A \textbf{Documentation} line at the end
is required (write \textbf{None} if you did not collaborate).}

\begin{questions}

\question \textbf{Reading a sweep.} A lab group sweeps the PHASER with all eight
elements active and a source at boresight. Working at $10.3\ \ghz$ with element
spacing $d = 14\ \text{mm}$, they record the trace below. The sweep step is
$2.8125\mydeg$ and the peak reads $-4.1$ dBFS.

\begin{center}
\begin{tabular}{r|rrrrrrrrrr}
$\theta$ (\mydeg) & $0.0$ & $2.8$ & $5.6$ & $8.4$ & $11.3$ & $14.1$ & $16.9$ & $19.7$ & $22.5$ & $25.3$ \\ \hline
Power (dBFS) & $-4.1$ & $-4.6$ & $-6.2$ & $-9.1$ & $-14.3$ & $-24.1$ & $-22.0$ & $-17.4$ & $-16.6$ & $-18.1$ \\
\end{tabular}
\end{center}

\begin{parts}

	\begin{minipage}[t]{0.58\linewidth}
		\part Using linear interpolation between the two samples that bracket the
		half-power level, find the half-power beamwidth. The trace is symmetric
		about $0\mydeg$.
		\begin{solution}[1.3in]
		Half power is $3$ dB below the peak, so the level is
		$-4.1 - 3.0 = -7.1$ dBFS. That falls between the samples at $5.6\mydeg$
		($-6.2$) and $8.4\mydeg$ ($-9.1$).
		\eq{ \theta_{-3} &= 5.6 + 2.8\lp\frac{7.1-6.2}{9.1-6.2}\rp = 5.6 + 2.8(0.31) = 6.5\mydeg }
		The trace is symmetric, so the beamwidth is twice this.
		\eq{ \theta_{\text{HP}} = 2(6.5) = 12.9\mydeg }
		Theory gives $13.2\mydeg$, so the sweep reads about three tenths of a
		degree narrow.
		\end{solution}
	\end{minipage}\hfill%
	\begin{minipage}[t]{0.37\linewidth}
		\raggedleft \vspace{12pt}
		$\theta_{\text{HP}}$ = \ansbox[1.3in]{$12.9\mydeg$}
	\end{minipage}

	\begin{minipage}[t]{0.58\linewidth}
		\part Find the first-null beamwidth as the sweep reports it, and compare it
		with $\text{FNBW} = 2\arcsin\lp\lambda/Nd\rp$.
		\begin{solution}[1.4in]
		The deepest sample between the main lobe and the sidelobe is $-24.1$ dBFS
		at $14.1\mydeg$, so the sweep reports
		\eq{ \text{FNBW}_{\text{meas}} = 2(14.1) = 28.2\mydeg }
		With $\lambda = 29.1\ \text{mm}$ and $Nd = 8(14) = 112\ \text{mm}$,
		\eq{ \text{FNBW} &= 2\arcsin\lp\frac{29.1}{112}\rp = 2\arcsin(0.2598) \\ &= 2(15.1\mydeg) = 30.1\mydeg }
		The measured value is $1.9\mydeg$ narrow because the true null at
		$15.1\mydeg$ falls between the samples at $14.1\mydeg$ and $16.9\mydeg$.
		The sweep can only report a null at an angle it actually visited.
		\end{solution}
	\end{minipage}\hfill%
	\begin{minipage}[t]{0.37\linewidth}
		\raggedleft \vspace{12pt}
		$\text{FNBW}_{\text{meas}}$ = \ansbox[1.3in]{$28.2\mydeg$}
		\\[10pt]
		$\text{FNBW}_{\text{calc}}$ = \ansbox[1.3in]{$30.1\mydeg$}
	\end{minipage}

	\begin{minipage}[t]{0.58\linewidth}
		\part Find the first sidelobe level relative to the peak, in dBc.
		\begin{solution}[1.1in]
		The tallest sample beyond the null is $-16.6$ dBFS at $22.5\mydeg$.
		Relative to the peak,
		\eq{ \text{SLL} = -16.6 - (-4.1) = -12.5\ \text{dBc} }
		which is within a few tenths of the $-12.8$ dB the array factor gives
		for eight uniform elements.
		\end{solution}
	\end{minipage}\hfill%
	\begin{minipage}[t]{0.37\linewidth}
		\raggedleft \vspace{12pt}
		SLL = \ansbox[1.3in]{$-12.5$ dBc}
	\end{minipage}

	\part No sample in the sweep reads below about $-24$ dBFS, which is $20$ dB
	below the peak, yet the noise floor on this kit is $23$ dB below the peak.
	Explain in one or two sentences why the deepest sample sits $3$ dB above the
	floor rather than on it.
		\begin{solution}[1.3in]
		The receiver reports signal power plus noise power. At the sample nearest
		the null the array factor has fallen to about $-23$ dB, which is the floor
		level itself, so the two contributions are equal and their sum is $3$ dB
		above either one. A measured null depth is therefore a statement about the
		dynamic range of the measurement rather than about the array.
		\end{solution}

\end{parts}

\newpage

\question \textbf{Predict before you measure.} The same array is reconfigured by
setting some element gains to zero. The spacing $d = 14\ \text{mm}$ and the
frequency $10.3\ \ghz$ do not change, so $\lambda = 29.1\ \text{mm}$ throughout.

\begin{parts}

	\begin{minipage}[t]{0.58\linewidth}
		\part Rx1, Rx2, Rx7 and Rx8 are set to zero, leaving Rx3 through Rx6
		active --- four adjacent elements. Predict the half-power beamwidth from
		$\theta_{\text{HP}} \approx 0.886\lambda/(Nd)$.
		\begin{solution}[1.2in]
		Four active elements span $Nd = 4(14) = 56\ \text{mm}$.
		\eq{ \theta_{\text{HP}} \approx \frac{0.886(29.1)}{56} = 0.4604\ \text{rad} = 26.4\mydeg }
		Solving the array factor exactly gives $27.4\mydeg$; the small-angle step
		inside the $0.886$ rule costs about a degree at this beamwidth.
		\end{solution}
	\end{minipage}\hfill%
	\begin{minipage}[t]{0.37\linewidth}
		\raggedleft \vspace{12pt}
		$\theta_{\text{HP}}$ = \ansbox[1.3in]{$\approx 26\text{-}27\mydeg$}
	\end{minipage}

	\begin{minipage}[t]{0.58\linewidth}
		\part Now only Rx3, Rx4 and Rx5 are left on --- three adjacent elements.
		Predict the half-power beamwidth and the first-null beamwidth.
		\begin{solution}[1.4in]
		Three adjacent elements span $Nd = 3(14) = 42\ \text{mm}$.
		\eq{ \theta_{\text{HP}} \approx \frac{0.886(29.1)}{42} = 0.6138\ \text{rad} = 35.2\mydeg }
		\eq{ \text{FNBW} &= 2\arcsin\lp\frac{29.1}{42}\rp = 2\arcsin(0.6929) \\ &= 2(43.8\mydeg) = 87.7\mydeg }
		The beam is wide enough that the $0.886$ rule is starting to under-predict;
		expect the sweep to read closer to $37\mydeg$.
		\end{solution}
	\end{minipage}\hfill%
	\begin{minipage}[t]{0.37\linewidth}
		\raggedleft \vspace{12pt}
		$\theta_{\text{HP}}$ = \ansbox[1.3in]{$\approx 35\mydeg$}
		\\[10pt]
		FNBW = \ansbox[1.3in]{$87.7\mydeg$}
	\end{minipage}

	\begin{minipage}[t]{0.58\linewidth}
		\part By how many dB should the peak of the three-element sweep sit below
		the peak of the full eight-element sweep, if every element gain that is
		still on remains at 100\%?
		\begin{solution}[1.1in]
		At boresight the element signals add in voltage, so the peak field is
		proportional to the number of active elements.
		\eq{ \Delta = 20\mylog{\frac{3}{8}} = 20\mylog{0.375} = -8.5\ \db }
		\end{solution}
	\end{minipage}\hfill%
	\begin{minipage}[t]{0.37\linewidth}
		\raggedleft \vspace{12pt}
		$\Delta$ = \ansbox[1.3in]{$-8.5\ \db$}
	\end{minipage}

	\part A student proposes reaching the same $35\mydeg$ beamwidth by leaving all
	eight elements on and turning every gain down to $37.5\%$. Explain in one or
	two sentences why that does not work.
		\begin{solution}[1.2in]
		Beamwidth is set by the active aperture $Nd$, not by how much gain the
		elements are given. Turning all eight gains down uniformly keeps the
		aperture at $112\ \text{mm}$, so the beam stays at about $13\mydeg$ and
		only the peak level drops. To widen the beam the aperture itself has to
		shrink.
		\end{solution}

\end{parts}

\newpage

\question \textbf{Reconciling a sidelobe.} A second group measures the first
sidelobe of the full eight-element array at $-11.5$ dBc, while the array factor
for $N = 8$ predicts $-12.8$ dB. Their noise floor sits at $-23$ dBc and the
sweep step is $2.8125\mydeg$.

\begin{parts}

	\begin{minipage}[t]{0.58\linewidth}
		\part The floor adds noise power to the sidelobe power the receiver
		reports. Treating the two as uncorrelated, compute the level the sweep
		should report for a true $-12.8$ dB sidelobe sitting on a $-23$ dBc
		floor.
		\begin{solution}[1.4in]
		Convert both to power ratios and add.
		\eq{ p_{\text{sig}} = 10^{-12.8/10} = 0.0525, \qquad p_{\text{n}} = 10^{-23/10} = 0.0050 }
		\eq{ p_{\text{tot}} = 0.0525 + 0.0050 = 0.0575 }
		\eq{ \text{SLL} = 10\mylog{0.0575} = -12.4\ \text{dBc} }
		The floor alone accounts for about $0.4$ dB of lift.
		\end{solution}
	\end{minipage}\hfill%
	\begin{minipage}[t]{0.37\linewidth}
		\raggedleft \vspace{12pt}
		SLL = \ansbox[1.3in]{$-12.4$ dBc}
	\end{minipage}

	\begin{minipage}[t]{0.58\linewidth}
		\part Their peak reading is also affected. The peak is $23$ dB above the
		floor by definition. By how much does the floor lift the peak reading, and
		which way does that push the reported sidelobe level in dBc?
		\begin{solution}[1.3in]
		Normalize the peak signal power to 1.
		\eq{ p_{\text{tot}} &= 1 + 0.0050 = 1.0050 \\ 10\mylog{1.0050} &= +0.02\ \db }
		The peak lifts by only about $0.02$ dB, which is negligible. Because the
		sidelobe lifts by $0.4$ dB and the peak barely moves, the ratio of the
		two moves \emph{up}, making the reported sidelobe look higher than the
		array factor predicts.
		\end{solution}
	\end{minipage}\hfill%
	\begin{minipage}[t]{0.37\linewidth}
		\raggedleft \vspace{12pt}
		peak lift = \ansbox[1.3in]{$+0.02\ \db$}
	\end{minipage}

	\part The floor explains $0.4$ dB of the $1.3$ dB discrepancy. Name two
	other contributions that plausibly make up the rest, and say for each whether
	it would move the reading up or down.
		\begin{solution}[1.5in]
		First, run-to-run noise on the sidelobe sample itself. The sidelobe is only
		about $10$ dB above the floor, so a single sweep reading scatters a few
		tenths of a decibel in either direction. Second, residual amplitude and
		phase errors across the eight elements after calibration. These break the
		uniform distribution slightly and raise the sidelobes, so that contribution
		moves the reading up. Grid straddle is a third contribution, but it moves
		the reading \emph{down} and by very little here, since the sidelobe peak at
		$21.9\mydeg$ is caught by the sample at $22.5\mydeg$ to within
		$0.05$ dB.
		\end{solution}

\end{parts}

\newpage

\question \textbf{The two-element case.} With only Rx4 and Rx5 active the array
has two elements at $d = 14\ \text{mm}$, and at $10.3\ \ghz$ the array factor is
$AF = \cosp{\psi/2}$ with $\psi = kd\sinp{\theta}$.

\begin{parts}

	\begin{minipage}[t]{0.58\linewidth}
		\part Find the half-power beamwidth exactly, without using the
		$0.886\lambda/(Nd)$ rule.
		\begin{solution}[1.4in]
		Set the array factor to $1/\sqrt{2}$.
		\eq{ \cosp{\psi/2} = \frac{1}{\sqrt{2}} \quad\Rightarrow\quad \frac{\psi}{2} = \frac{\pi}{4} \quad\Rightarrow\quad \psi = \frac{\pi}{2} }
		With $\psi = kd\sinp{\theta}$ and $kd = 2\pi d/\lambda$,
		\eq{ \sinp{\theta} = \frac{\pi/2}{2\pi d/\lambda} = \frac{\lambda}{4d} = \frac{29.1}{4(14)} = 0.5196 }
		\eq{ \theta = 31.3\mydeg \quad\Rightarrow\quad \theta_{\text{HP}} = 62.6\mydeg }
		\end{solution}
	\end{minipage}\hfill%
	\begin{minipage}[t]{0.37\linewidth}
		\raggedleft \vspace{12pt}
		$\theta_{\text{HP}}$ = \ansbox[1.3in]{$62.6\mydeg$}
	\end{minipage}

	\begin{minipage}[t]{0.58\linewidth}
		\part Show that the $0.886\lambda/(Nd)$ rule fails here, and state the
		percentage error it makes.
		\begin{solution}[1.2in]
		\eq{ \theta_{\text{HP}} \approx \frac{0.886(29.1)}{28} = 0.9208\ \text{rad} = 52.8\mydeg }
		\eq{ \text{error} = \frac{52.8 - 62.6}{62.6} = -15.7\% }
		The rule assumes the half-power point is at a small enough angle that
		$\sinp{\theta} \approx \theta$. At $31\mydeg$ that is no longer true.
		\end{solution}
	\end{minipage}\hfill%
	\begin{minipage}[t]{0.37\linewidth}
		\raggedleft \vspace{12pt}
		error = \ansbox[1.3in]{$-15.7\%$}
	\end{minipage}

	\begin{minipage}[t]{0.58\linewidth}
		\part Evaluate $\lambda/2d$ and use it to explain why the lab table lists
		the first-null beamwidth of this array as $180\mydeg$.
		\begin{solution}[1.3in]
		The first null of $\cosp{\psi/2}$ needs $\psi = \pi$, which requires
		\eq{ \sinp{\theta} = \frac{\lambda}{2d} = \frac{29.1}{2(14)} = 1.039 }
		No real angle has a sine greater than 1, so the null lies outside visible
		space and no null appears anywhere in the $\pm 90\mydeg$ scan. With no
		null to measure, the beamwidth between first nulls is quoted as the full
		$180\mydeg$ scan range by convention.
		\end{solution}
	\end{minipage}\hfill%
	\begin{minipage}[t]{0.37\linewidth}
		\raggedleft \vspace{12pt}
		$\lambda/2d$ = \ansbox[1.3in]{$1.039$}
	\end{minipage}

	\begin{minipage}[t]{0.58\linewidth}
		\part How far below its own peak does this two-element pattern fall at
		$\theta = 90\mydeg$? Compare that with the $23$ dB of dynamic range the
		sweep has above the noise floor at eight elements, and say what the trace
		will look like near the edges of the scan.
		\begin{solution}[1.5in]
		At $\theta = 90\mydeg$, $\sinp{\theta} = 1$ and
		$\psi = 2\pi(14)/29.1 = 3.023$ rad, so
		\eq{ \vert AF \vert &= \cosp{1.511} = 0.0597 \\ 20\mylog{0.0597} &= -24.5\ \db }
		The two-element peak is itself $12$ dB below the eight-element peak, so
		it clears the floor by only about $11$ dB. The pattern needs
		$24.5$ dB of range to reach endfire and the measurement has $11$ dB,
		so the trace flattens into the noise well before $\pm 90\mydeg$ and the
		outer part of the roll-off is not visible.
		\end{solution}
	\end{minipage}\hfill%
	\begin{minipage}[t]{0.37\linewidth}
		\raggedleft \vspace{12pt}
		drop at $90\mydeg$ = \ansbox[1.3in]{$-24.5\ \db$}
	\end{minipage}

\end{parts}

\newpage

\question \textbf{Design judgment.} Your section is asked to demonstrate a
$20\mydeg$ beam on the PHASER at $10.3\ \ghz$, using only the Element Gains
controls.

\begin{parts}

	\begin{minipage}[t]{0.58\linewidth}
		\part How many adjacent elements must be active to give a half-power
		beamwidth closest to $20\mydeg$? Give the aperture and the resulting
		beamwidth.
		\begin{solution}[1.4in]
		Solve the beamwidth rule for the aperture.
		\eq{ Nd = \frac{0.886\lambda}{\theta_{\text{HP}}} = \frac{0.886(29.1)}{20\mydeg(\pi/180)} = \frac{25.78}{0.3491} = 73.8\ \text{mm} }
		\eq{ N = \frac{73.8}{14} = 5.3 }
		Element counts are integers, so take $N = 5$: $Nd = 70\ \text{mm}$ and
		\eq{ \theta_{\text{HP}} \approx \frac{0.886(29.1)}{70} = 0.3683\ \text{rad} = 21.1\mydeg }
		Six elements would give $17.6\mydeg$, which is further from the target.
		\end{solution}
	\end{minipage}\hfill%
	\begin{minipage}[t]{0.37\linewidth}
		\raggedleft \vspace{12pt}
		$N$ = \ansbox[1.3in]{$5$ elements}
		\\[10pt]
		$\theta_{\text{HP}}$ = \ansbox[1.3in]{$21.1\mydeg$}
	\end{minipage}

	\begin{minipage}[t]{0.58\linewidth}
		\part What does this cost in peak received level compared with the full
		eight-element array, and how much dynamic range is left above the noise
		floor?
		\begin{solution}[1.3in]
		\eq{ \Delta = 20\mylog{\frac{5}{8}} = 20\mylog{0.625} = -4.1\ \db }
		The floor does not move, so the range above it drops from about $23$ dB
		to
		\eq{ 23 - 4.1 = 18.9\ \db }
		The first sidelobe of a five-element uniform array is near $-12$ dB,
		which still clears the floor comfortably.
		\end{solution}
	\end{minipage}\hfill%
	\begin{minipage}[t]{0.37\linewidth}
		\raggedleft \vspace{12pt}
		$\Delta$ = \ansbox[1.3in]{$-4.1\ \db$}
		\\[10pt]
		range = \ansbox[1.3in]{$18.9\ \db$}
	\end{minipage}

	\part A classmate suggests instead applying a Hann taper to all eight
	elements, since that also widens the beam to about $19.5\mydeg$. State one
	advantage and one disadvantage of the taper against switching elements off,
	and say which you would choose for a demonstration whose purpose is to show
	that beamwidth depends on aperture.
		\begin{solution}[1.6in]
		The taper keeps the full $112\ \text{mm}$ aperture and pushes the
		sidelobes below the noise floor, so the trace is cleaner and the main lobe
		is the only thing on the plot. Its disadvantage for this purpose is that
		it widens the beam by changing the shape of the amplitude distribution,
		not by changing the aperture, so it demonstrates the opposite of the
		intended point. It also costs about $4.7$ dB of peak, which is no better
		than the five-element case. For a demonstration about aperture, switch
		elements off: the beamwidth and the aperture then move together and the
		claim is visible in one control.
		\end{solution}

\end{parts}

%%%%%%%%%%%%%%%%%%%%%%%%%%%%%%%%%%%%%%%%%%%%%%%%%%%
\newpage
\fullwidth{\textbf{Part B -- Learning Objective 3.6: } I can distinguish between array factor
and true antenna pattern and account for element pattern effects.}
\question \textbf{Reading dBc off a Tracking-tab trace.} A cadet walks an HB100
around the PHASER on a $1\ \m$ arc with the beam held at broadside and records the
Signal vs Time trace. The three lobe peaks read $-18.4\ \text{dBFS}$ (main lobe),
$-31.2\ \text{dBFS}$ (left first sidelobe), and $-29.9\ \text{dBFS}$ (right first
sidelobe).

\begin{parts}

	\begin{minipage}[t]{0.58\linewidth}
		\part Convert both sidelobe readings to dBc.
		\begin{solution}[1.1in]
		A dBc value is the lobe peak referred to the main-lobe peak, so subtract:
		\eq{ L_{\text{left}} &= -31.2 - (-18.4) = -12.8\ \text{dBc} \\
		     L_{\text{right}} &= -29.9 - (-18.4) = -11.5\ \text{dBc} }
		\end{solution}
	\end{minipage}\hfill%
	\begin{minipage}[t]{0.37\linewidth}
		\raggedleft \vspace{12pt}
		$L_{\text{left}}$ = \ansbox[1.25in]{$-12.8$ dBc}
		\\[12pt]
		$L_{\text{right}}$ = \ansbox[1.25in]{$-11.5$ dBc}
	\end{minipage}

	\part The cadet's partner raises \textbf{Rx Gain} by $6\ \db$ and repeats the
	walk. State what happens to the three dBFS readings and what happens to the
	two dBc values, and explain why the two answers differ.
		\begin{solution}[1.3in]
		All three dBFS readings rise by about $6\ \db$ as well, because the gain change
		multiplies every received sample by the same factor. The dBc values do not
		change, because each one is a difference between two readings taken through
		the same gain setting and the common term cancels in the subtraction. This
		is why a pattern is reported as a ratio to its own peak: the result does not
		depend on the receiver's absolute scale, which is not calibrated here anyway.
		\end{solution}

	\part A uniform 8-element array has its first sidelobes at $-13$ dBc. Comment
	on the pair measured above: is this array uniformly illuminated, and what
	accounts for the $1.3\ \db$ spread between the two sides?
		\begin{solution}[1.4in]
		Both readings straddle $-13$ dBc within about a decibel and a half, which is
		consistent with uniform illumination measured in a room. A tapered
		illumination would put the sidelobes at least $17\ \db$ down, and a badly
		calibrated array would show a much larger imbalance. The $1.3\ \db$ spread
		between left and right comes from the measurement rather than the array:
		classroom multipath adds roughly $\pm 1\ \db$ of ripple, and the walking speed
		and radius were not identical on the two halves of the arc. The array factor
		of a symmetric uniform array is symmetric at broadside, so the asymmetry has
		no home in the physics of the array itself.
		\end{solution}

\end{parts}

\newpage

\question \textbf{Why the time axis is not an angle axis.} The same trace runs from
$t = 0.00\ \text{s}$, when the source was at $-90\mydeg$, to $t = 7.40\ \text{s}$,
when it reached $+90\mydeg$. The main-lobe peak lands at $t = 3.10\ \text{s}$. The
beam was commanded to broadside.

\begin{parts}

	\part Explain why the lobe \emph{amplitudes} on this trace are usable while the
	lobe \emph{angles} are not.
		\begin{solution}[1.5in]
		The amplitude recorded at any instant is the pattern evaluated at whatever
		angle the source occupied at that instant, and that value does not depend on
		how fast the source was moving. Walking slowly through a lobe spends more
		samples on it and draws it wider on the plot, but the height each sample
		reaches is unchanged. The angle, in contrast, is never recorded at all. The
		x-axis stores elapsed time, and recovering angle from time requires knowing
		the walking speed at every instant, which nothing in the measurement captures.
		\end{solution}

	\begin{minipage}[t]{0.58\linewidth}
		\part Assume a constant walking speed and convert the main-lobe time to an
		angle. Compare it against the commanded beam angle and state the error.
		\begin{solution}[1.2in]
		A constant speed maps time linearly onto the $180\mydeg$ swept:
		\eq{ \theta &= -90\mydeg + 180\mydeg \lp \frac{3.10}{7.40} \rp \\
		            &= -90\mydeg + 75.4\mydeg = -14.6\mydeg }
		The beam was commanded to $0\mydeg$, so the linear map is in error by
		$14.6\mydeg$.
		\end{solution}
	\end{minipage}\hfill%
	\begin{minipage}[t]{0.37\linewidth}
		\raggedleft \vspace{12pt}
		$\theta$ = \ansbox[1.25in]{$-14.6\mydeg$}
		\\[12pt]
		error = \ansbox[1.25in]{$14.6\mydeg$}
	\end{minipage}

	\begin{minipage}[t]{0.58\linewidth}
		\part The true half-power beamwidth is $13.1\mydeg$. If the hand moves as
		much as $30\%$ faster or slower than its average while crossing the main
		lobe, give the range of beamwidths a constant-speed conversion would report.
		\begin{solution}[1.2in]
		A hand moving $30\%$ fast spends $1/1.30$ of the expected time in the lobe,
		and a constant-speed map converts that shortened interval into a narrower
		angle. A hand moving $30\%$ slow does the opposite:
		\eq{ \theta_{\text{fast}} &= \frac{13.1\mydeg}{1.30} = 10.1\mydeg \\
		     \theta_{\text{slow}} &= \frac{13.1\mydeg}{0.70} = 18.7\mydeg }
		The reported beamwidth spans nearly a factor of two on a pattern that never
		changed.
		\end{solution}
	\end{minipage}\hfill%
	\begin{minipage}[t]{0.37\linewidth}
		\raggedleft \vspace{12pt}
		range = \ansbox[1.4in]{$10.1\mydeg$ to $18.7\mydeg$}
	\end{minipage}

	\part Name the piece of range hardware that removes this error and state what
	it records that the hand walk does not.
		\begin{solution}[1.2in]
		A turntable with a shaft encoder. The turntable carries the antenna through a
		commanded rotation, and the encoder tags every recorded sample with the
		angular position it was taken at, typically to a small fraction of a degree.
		The measurement then stores angle directly instead of inferring it from time,
		so no assumption about rotation speed is needed and a non-uniform rotation
		costs nothing.
		\end{solution}

\end{parts}

\newpage

\question \textbf{Scan loss and sidelobe asymmetry.} The projected-aperture rule
puts the peak power gain of a steered array at $\cos\theta_0$ relative to broadside.

\begin{parts}

	\begin{minipage}[t]{0.58\linewidth}
		\part Compute the predicted peak drop at steer angles of $30\mydeg$,
		$45\mydeg$, and $60\mydeg$.
		\begin{solution}[1.3in]
		The rule applies to power, so the drop in dB is $10\mylog{\cos\theta_0}$:
		\eq{ 30\mydeg&: \ 10\mylog{0.866} = -0.6\ \db \\
		     45\mydeg&: \ 10\mylog{0.707} = -1.5\ \db \\
		     60\mydeg&: \ 10\mylog{0.500} = -3.0\ \db }
		\end{solution}
	\end{minipage}\hfill%
	\begin{minipage}[t]{0.37\linewidth}
		\raggedleft \vspace{12pt}
		$30\mydeg$ = \ansbox[1.1in]{$-0.6\ \db$}
		\\[10pt]
		$45\mydeg$ = \ansbox[1.1in]{$-1.5\ \db$}
		\\[10pt]
		$60\mydeg$ = \ansbox[1.1in]{$-3.0\ \db$}
	\end{minipage}

	\begin{minipage}[t]{0.58\linewidth}
		\part A broadside walk peaks at $-18.4\ \text{dBFS}$ and a $30\mydeg$ walk
		peaks at $-19.3\ \text{dBFS}$. Give the measured scan loss, and state whether
		one pair of walks resolves it given $\pm 0.5\ \db$ of run-to-run scatter.
		\begin{solution}[1.3in]
		\eq{ \Delta = -19.3 - (-18.4) = -0.9\ \db }
		The measured $0.9\ \db$ sits $0.3\ \db$ from the predicted value of $0.6\ \db$ and well inside
		the $\pm 0.5\ \db$ scatter of a hand walk. One pair of runs is therefore
		consistent with the prediction but does not resolve it: the same pair of
		numbers would also be consistent with a scan loss anywhere between about
		$0.4$ and $1.4\ \db$ of scan loss. Averaging several walks at each steer angle narrows the
		scatter enough to make the comparison worth reporting.
		\end{solution}
	\end{minipage}\hfill%
	\begin{minipage}[t]{0.37\linewidth}
		\raggedleft \vspace{12pt}
		$\Delta$ = \ansbox[1.25in]{$-0.9\ \db$}
	\end{minipage}

	\part The measured $30\mydeg$ trace shows first sidelobes of $-12.2$ dBc on the
	low-angle side and $-15.2$ dBc on the high-angle side. Explain why they are
	unequal, given that the array factor alone would make them equal.
		\begin{solution}[1.4in]
		The array factor of a uniformly excited array is a function of
		$\sin\theta - \sin\theta_0$, so its sidelobes are equal in that variable and
		the two first sidelobes of a steered beam carry the same array-factor height.
		The radiated pattern is the array factor multiplied by the element factor,
		and the element factor falls off monotonically away from broadside. The
		high-angle sidelobe near $+59\mydeg$ therefore sits where the element pattern
		is much weaker than at the low-angle sidelobe near $+8\mydeg$, and the
		element factor pulls it down by the difference. Asymmetry in a steered
		pattern is evidence of the element factor, not of a fault in the array.
		\end{solution}

	\begin{minipage}[t]{0.58\linewidth}
		\part Real patch elements follow a power pattern nearer
		$\cos^{1.4}\theta$ than $\cos\theta$. Recompute the $60\mydeg$ scan loss and
		state how much worse it is than the ideal-element bound.
		\begin{solution}[1.2in]
		\eq{ 10\mylog{\cos^{1.4} 60\mydeg} &= 14\mylog{0.500} \\ &= -4.2\ \db }
		That is $1.2\ \db$ worse than the $-3.0\ \db$ ideal-element bound. The
		$\cos\theta$ rule is a best case set by the projected aperture alone; a real
		element's own pattern adds to it, and a link budget that scans to $60\mydeg$
		should carry the larger number.
		\end{solution}
	\end{minipage}\hfill%
	\begin{minipage}[t]{0.37\linewidth}
		\raggedleft \vspace{12pt}
		loss = \ansbox[1.25in]{$-4.2\ \db$}
	\end{minipage}

\end{parts}

\newpage

\question \textbf{How far out is far enough.} The PHASER's 8 elements are spaced
$14\ \text{mm}$ apart, so the aperture spans $D = 98\ \text{mm}$ between the outer
element centers. The HB100 runs at $10.525\ \ghz$.

\begin{parts}

	\begin{minipage}[t]{0.58\linewidth}
		\part Find the far-field distance for this aperture and state whether the
		$1\ \m$ arc used in lab is far enough.
		\begin{solution}[1.3in]
		\eq{ \lambda &= \frac{3\times 10^8}{10.525\times 10^9} = 0.0285\ \m \\
		     r &\ge \frac{2D^2}{\lambda} = \frac{2(0.098)^2}{0.0285} = 0.67\ \m }
		The $1\ \m$ arc is comfortably outside the $0.67\ \m$ boundary, so the pattern has settled
		into its far-field shape everywhere along the walk.
		\end{solution}
	\end{minipage}\hfill%
	\begin{minipage}[t]{0.37\linewidth}
		\raggedleft \vspace{12pt}
		$r$ = \ansbox[1.25in]{$0.67\ \m$}
	\end{minipage}

	\begin{minipage}[t]{0.58\linewidth}
		\part A Ku-band panel with a $0.200\ \m$ aperture is to be measured at
		$15.0\ \ghz$. Find its far-field distance.
		\begin{solution}[1.2in]
		\eq{ \lambda &= \frac{3\times 10^8}{15.0\times 10^9} = 0.0200\ \m \\
		     r &\ge \frac{2(0.200)^2}{0.0200} = 4.00\ \m }
		Doubling the aperture and halving the wavelength pushes the requirement out
		by a factor of eight.
		\end{solution}
	\end{minipage}\hfill%
	\begin{minipage}[t]{0.37\linewidth}
		\raggedleft \vspace{12pt}
		$r$ = \ansbox[1.25in]{$4.00\ \m$}
	\end{minipage}

	\begin{minipage}[t]{0.58\linewidth}
		\part The Ku-band panel is measured on the same $1.0\ \m$ classroom arc anyway.
		Find the phase error across the aperture edge and compare it against the
		$22.5\mydeg$ tolerance behind the far-field criterion.
		\begin{solution}[1.3in]
		The extra path to the aperture edge relative to its center is
		$\Delta \approx D^2/8r$, and the phase error is $360\mydeg\Delta/\lambda$:
		\eq{ \Delta &= \frac{(0.200)^2}{8(1.0)} = 0.005\ \m \\
		     \Delta\phi &= \frac{360\mydeg (0.005)}{0.0200} = 90\mydeg }
		That is four times the $22.5\mydeg$ the criterion allows.
		\end{solution}
	\end{minipage}\hfill%
	\begin{minipage}[t]{0.37\linewidth}
		\raggedleft \vspace{12pt}
		$\Delta\phi$ = \ansbox[1.25in]{$90\mydeg$}
	\end{minipage}

	\part Describe what a $90\mydeg$ edge phase error does to the measured pattern,
	and say which pattern feature it corrupts worst.
		\begin{solution}[1.4in]
		The edges of the aperture arrive out of phase with the center, so the
		measurement no longer sees a plane wave and the aperture is effectively
		defocused. The main lobe broadens and its peak drops, but the damage is worst
		in the nulls and the sidelobes: the cancellation that produces a deep null
		depends on precise phase alignment across the aperture, so the nulls fill in
		and the sidelobes read several decibels high. A sidelobe level measured this
		way is pessimistic, and a null depth measured this way is meaningless.
		\end{solution}

\end{parts}

\newpage

\question \textbf{Classroom against chamber.} The lab trace carries about
$\pm 1\ \db$ of ripple from room reflections adding to the direct path.

\begin{parts}

	\part List which of the four hand-measurement artifacts from the lab an
	anechoic chamber removes, and which one it does not remove by itself.
		\begin{solution}[1.5in]
		A chamber removes the reflection artifact directly: absorber on every surface
		leaves only the direct path, so the ripple disappears. A chamber's fixtures
		also remove the angle error and the amplitude instability, because the
		antenna sits on an encoded turntable and the source is bolted at a fixed
		range rather than carried. The far-field condition is not removed by the
		absorber at all; it is a geometric requirement on the separation distance,
		and a chamber satisfies it only if the chamber is long enough for the
		aperture and frequency being measured. A large antenna in a short chamber has
		the same problem as a large antenna in a classroom.
		\end{solution}

	\begin{minipage}[t]{0.58\linewidth}
		\part A single reflected path of relative amplitude $\rho$ adds to a direct
		path of amplitude $1$. Find the $\rho$ that produces a $+1\ \db$ peak in the
		received amplitude, and express it in dB.
		\begin{solution}[1.3in]
		The two paths add in phase at the ripple peak:
		\eq{ 20\mylog{1+\rho} &= 1.0\ \db \\
		     1+\rho &= 10^{0.05} = 1.122 \\
		     \rho &= 0.122 \ \rightarrow \ -18.3\ \db }
		A reflection only $18\ \db$ below the direct path is enough to move the
		reading by a decibel, which is why room reflections matter at levels that
		sound negligible.
		\end{solution}
	\end{minipage}\hfill%
	\begin{minipage}[t]{0.37\linewidth}
		\raggedleft \vspace{12pt}
		$\rho$ = \ansbox[1.25in]{$0.122 = -18.3\ \db$}
	\end{minipage}

	\begin{minipage}[t]{0.58\linewidth}
		\part The main-lobe peak and the sidelobe peak each carry their own
		$\pm 1\ \db$ of ripple. Give the worst-case uncertainty on a sidelobe level
		reported in dBc, and give the range a true $-13$ dBc sidelobe could read.
		\begin{solution}[1.2in]
		A dBc value is a difference of two readings, and in the worst case the two
		errors have opposite signs and add:
		\eq{ \pm 1\ \db + \pm 1\ \db = \pm 2\ \db }
		A true $-13$ dBc sidelobe can therefore be reported anywhere from $-11$ to
		$-15$ dBc without anything being wrong with the array.
		\end{solution}
	\end{minipage}\hfill%
	\begin{minipage}[t]{0.37\linewidth}
		\raggedleft \vspace{12pt}
		uncertainty = \ansbox[1.1in]{$\pm 2\ \db$}
		\\[12pt]
		range = \ansbox[1.25in]{$-11$ to $-15$ dBc}
	\end{minipage}

	\part A customer requires written proof that this array's sidelobes are below
	$-25$ dBc. State whether the classroom measurement can support that claim, and
	give the two reasons that decide it.
		\begin{solution}[1.4in]
		It cannot. The first reason is the uncertainty computed above: a $\pm 2\ \db$
		window on every dBc reading cannot certify a level that must sit below a
		specific number, because a measured $-26$ dBc could be a true $-24$ dBc. The
		second reason is dynamic range. The electrical sweep's noise floor sits about
		$23\ \db$ below the peak, so a $-25$ dBc sidelobe is buried in the floor and
		there is nothing to read. Both problems are what a chamber and a calibrated
		receiver are for, and the claim needs a range measurement.
		\end{solution}

\end{parts}

\vspace{1.5em}
\noindent\textbf{Documentation:} \rule[-2pt]{0.6\linewidth}{0.4pt}

\end{questions}
